\chapter{Overview}

This blueprint records the mathematical dependency structure of the
\href{https://github.com/kdwong/surreal}{surreal Lean repository}.  It follows the declarations
and proofs in that development rather than presenting a separate textbook construction.

Every game considered here has finitely many options at each node and finite birthday.
Accordingly, the quotient constructed here is the ring of
\emph{short} surreal numbers, denoted in this document by \(\NoShort\).  This is the dyadic
part of Conway's full class of surreal numbers.  No claim about proper-class-sized or
transfinite birthdays is made.

The principal dependency path is
\[
\begin{aligned}
  &\text{finite games and birthdays}
  \longrightarrow \text{recursive game order}
  \longrightarrow \text{surreal totality} \\
  &\quad\longrightarrow \text{ordered additive quotient},\\[2mm]
  &\text{recursive product}
  \longrightarrow \text{Conway A/B/C induction}
  \longrightarrow \text{product closure and congruence}\\
  &\quad\longrightarrow \text{commutative ring}
  \longrightarrow \text{positive-product theorem}
  \longrightarrow \text{linear ordered commutative ring}.
\end{aligned}
\]

On the web version, each formalized item links to its Lean declaration.  The dependency
graph uses dashed edges for dependencies occurring in a statement or construction and solid
edges for proof dependencies.  The graph retains every declared edge, including transitively
redundant ones, so the complete audited dependency path remains visible.

All nodes marked as formalized were checked against the current local formal development.  A
green node means that the corresponding theorem has been verified; the prose proof explains
its mathematical argument and its place in the dependency path.
